\section{Banach--Stone theorem}

\begin{definition}[Faces of the unit sphere]
\label{def:banach-stone-face}
\lean{BanachStoneTheorem.face}
\leanok
For a compact space $X$, a point $x\in X$, and $\varepsilon\in\mathbb R$, define
\[
  F(x,\varepsilon)=\{f\in C(X,\mathbb R):\|f\|=1\text{ and }f(x)=\varepsilon\}.
\]
\end{definition}

\begin{lemma}[The norm is attained]
\label{lem:banach-stone-norm-attained}
\lean{BanachStoneTheorem.exists_abs_apply_eq_norm}
\leanok
If $X$ is compact and nonempty and $f\in C(X,\mathbb R)$, then there is $x\in X$ such that
$|f(x)|=\|f\|$.
\end{lemma}

\begin{lemma}[Faces are convex]
\label{lem:banach-stone-face-convex}
\lean{BanachStoneTheorem.face_convex}
\leanok
\uses{def:banach-stone-face}
If $|\varepsilon|=1$, then $F(x,\varepsilon)$ is convex.
\end{lemma}

\begin{lemma}[Finite common-sign property]
\label{lem:banach-stone-finite-common-sign}
\lean{BanachStoneTheorem.finite_common_sign}
\leanok
\uses{lem:banach-stone-norm-attained}
Let $A$ be a convex subset of the unit sphere of $C(X,\mathbb R)$.  Every finite
subset of $A$ has a point $x\in X$ and a common sign $\varepsilon\in\{-1,1\}$ at
which all its functions take the value $\varepsilon$.
\end{lemma}

\begin{lemma}[Convex subsets lie in a face]
\label{lem:banach-stone-convex-subset-face}
\lean{BanachStoneTheorem.convex_subset_face}
\leanok
\uses{def:banach-stone-face,lem:banach-stone-finite-common-sign}
Every convex subset of the unit sphere of $C(X,\mathbb R)$ is contained in a
face $F(x,\varepsilon)$ with $\varepsilon\in\{-1,1\}$.
\end{lemma}

\begin{lemma}[Separation by a function in a face]
\label{lem:banach-stone-separator}
\lean{BanachStoneTheorem.separator_mem_face}
\leanok
\uses{def:banach-stone-face}
If $x\ne y$, then for either sign there is a continuous function in the face at
$x$ taking the opposite sign at $y$.
\end{lemma}

\begin{lemma}[A face determines its point and sign]
\label{lem:banach-stone-face-rigidity}
\lean{BanachStoneTheorem.face_subset_face}
\leanok
\uses{lem:banach-stone-separator}
If one signed face is contained in another, then the underlying points and signs
are equal.
\end{lemma}

\begin{lemma}[Faces are maximal convex subsets]
\label{lem:banach-stone-face-maximal}
\lean{BanachStoneTheorem.face_maximal}
\leanok
\uses{lem:banach-stone-convex-subset-face,lem:banach-stone-face-rigidity}
Each signed face is a maximal convex subset of the unit sphere.
\end{lemma}

\begin{lemma}[Linear isometries send faces to faces]
\label{lem:banach-stone-isometry-faces}
\lean{BanachStoneTheorem.exists_image_face_eq_face}
\leanok
\uses{lem:banach-stone-face-convex,lem:banach-stone-convex-subset-face,lem:banach-stone-face-maximal}
A linear isometry between spaces of real-valued continuous functions maps each
signed face onto a signed face.
\end{lemma}

\begin{definition}[The induced point map]
\label{def:banach-stone-point-map}
\lean{BanachStoneTheorem.pointMap}
\leanok
\uses{lem:banach-stone-isometry-faces}
The image-of-faces correspondence defines a map between the underlying compact
Hausdorff spaces.
\end{definition}

\begin{lemma}[Pointwise form of the isometry]
\label{lem:banach-stone-map-apply}
\lean{BanachStoneTheorem.map_apply_eq}
\leanok
\uses{def:banach-stone-point-map,lem:banach-stone-isometry-faces}
For every $f$ and $y$, evaluation of the isometry at $y$ is, up to a sign,
evaluation of $f$ at the induced point.
\end{lemma}

\begin{lemma}[Continuity of the induced point map]
\label{lem:banach-stone-point-map-continuous}
\lean{BanachStoneTheorem.continuous_pointMap}
\leanok
\uses{lem:banach-stone-map-apply}
The induced point map is continuous.
\end{lemma}

\begin{lemma}[The point maps are inverse]
\label{lem:banach-stone-point-map-inverse}
\lean{BanachStoneTheorem.pointMap_symm_apply}
\leanok
\uses{lem:banach-stone-map-apply}
The point map associated with an isometry and the point map associated with its
inverse are mutually inverse.
\end{lemma}

\begin{definition}[The induced equivalence]
\label{def:banach-stone-point-equiv}
\lean{BanachStoneTheorem.pointEquiv}
\leanok
\uses{def:banach-stone-point-map,lem:banach-stone-point-map-inverse}
The two inverse point maps define an equivalence between the underlying spaces.
\end{definition}

\begin{theorem}[Banach--Stone]
\label{thm:banach-stone}
\lean{BanachStoneTheorem.MainTheorem}
\leanok
\uses{def:banach-stone-point-equiv,lem:banach-stone-point-map-continuous}
Let $X$ and $Y$ be compact Hausdorff spaces.  If $C(X,\mathbb R)$ and
$C(Y,\mathbb R)$ are linearly isometric, then $X$ and $Y$ are homeomorphic.
\end{theorem}
